% TOP_PROBLEM: {"id": 30005443, "title": "Zero–One Law for Directional Transience", "category_id": 19, "category": {"id": 19, "name": "probability", "display_name": "Probability", "description": "Problems involving probability theory, stochastic processes, and random structures.", "slug": "probability", "order_index": 19, "created_at": "2026-07-31T15:26:25.671Z"}, "status": "open"}
% FIRST_POSTED: null
% ENTRY_KIND: statement_only
% !TeX program = lualatex
\documentclass[11pt]{article}
\usepackage[margin=25mm]{geometry}
\usepackage{fontspec}
\setmainfont{Latin Modern Roman}
\usepackage{amsmath,amssymb,mathtools}
\usepackage{unicode-math}
\setmathfont{Latin Modern Math}
\usepackage{xurl,hyperref}
\hypersetup{colorlinks=true,urlcolor=blue}
\setcounter{secnumdepth}{0}
\setlength{\parindent}{0pt}
\setlength{\parskip}{5pt}
\setlength{\emergencystretch}{3em}
\newcommand{\R}{\mathbb R}
\newcommand{\N}{\mathbb N}
\newcommand{\Z}{\mathbb Z}
\newcommand{\Q}{\mathbb Q}
\newcommand{\C}{\mathbb C}
\newcommand{\D}{\mathbb D}
\newcommand{\F}{\mathbb F}
\newcommand{\sref}[2]{\hyperlink{source:#1:#2}{[S#2]}}
\newcommand{\cataloguescope}[1]{#1}
\begin{document}
% UnsolvedMath ID 30005443; source catalogue code OWR-12697706-001
\setcounter{section}{30005442}
\section{Zero–One Law for Directional Transience}\label{30005443}
{\small UnsolvedMath ID 30005443 · Probability}
\subsection{Definitions and mathematical statement}

\paragraph{Shared notation.}$\mathbb N$, $\mathbb Z$, $\mathbb Q$, $\mathbb R$, and $\mathbb C$ denote the natural numbers, integers, rationals, reals, and complex numbers, respectively. All additional parameters and hypotheses are specified in the question.


\paragraph{Statement recorded in the supplied source.}
For a random walk in a random environment, let $A_+$ be the event of directional transience in a fixed direction. Must $\mathbb P(A_+)\in\{0,1\}$?
\par\smallskip\textit{Formulation references:} \sref{30005443}{1}, \sref{30005443}{2}.
\subsection{Short English statement}
For a random walk in a random environment, let $A_+$ be the event of directional transience in a fixed direction. Must $\mathbb P(A_+)\in\{0,1\}$?
\subsection{Sources}
\begin{itemize}
\item[\hypertarget{source:30005443:1}{S1}] ULAM.AI and the UnsolvedMath Contributors, UnsolvedMath: A Curated Collection of Open Mathematics Problems, record 30005443 (OWR-12697706-001). Catalogue attribution under CC BY 4.0.
\url{https://huggingface.co/datasets/ulamai/UnsolvedMath}.
\item[\hypertarget{source:30005443:2}{S2}] Original problem formulation and mathematical context.
\url{https://doi.org/10.4171/owr/2023/10}.
\end{itemize}
\label{end:30005443}
\end{document}
